Structured pruning removes whole attention heads and feed-forward channel groups, producing smaller dense language models that run efficiently on standard hardware. Yet training-free methods degrade sharply at moderate-to-high pruning ratios. We attribute this failure to the one-shot paradigm: a surrogate built once at the dense model is extrapolated across the entire removal budget and becomes progressively less reliable as the model leaves its construction point. We introduce \method, a calibration-only anchored continuation that constructs a nested, cost-indexed pruning trajectory instead of selecting the target mask in one step. The dense teacher and its Fisher metric remain fixed as the anchor, while marginal ablation perturbations are refreshed at successive partially pruned states. The resulting fixed-metric surrogate incorporates a path-dependent drift term that accounts for damage accumulated during earlier removals; an interaction-guided partial-refresh strategy limits the additional estimation cost. Across seven language models from five architecture families, \method outperforms its one-shot counterpart in all 28 evaluated model--ratio settings. At a $50\%$ pruning ratio, it achieves the lowest WikiText-2 perplexity among the evaluated selection-only methods on all seven models, reducing perplexity by up to $16.8\times$ over one-shot without weight updates or additional deployment cost. The gains extend to downstream accuracy, while white-box mask analysis shows that continuation identifies removable attention heads that dense-point selection leaves protected.
